\chapter{Conclusion}
\chaptermark{Conclusion}
\label{chap:conclusion}

%% \begin{quote}
%% Chapter abstract goes here.
%% \end{quote}

%% \section{Introduction}
%% Introduce my dissertation topic.

\section{Future directions}

The emotional arcs of movies could be considered as a feature driving once controversial movies towards normalization over time, a closer examination of the trend presented by \cite{amendola2015evolving}.

Narratives are not unique in their explanation of causal links between events, and often the ``adjacent narratives'' are in direct competition.
We saw in Chapter 1 that the the disambiguation of competing event chains is an active area of NLP research.
This is identified as one factor contributing to Information overload \cite{orman2015information} via the internet, and participating in a collective cognitive denial of service attack \citep{king2016chinese}.
In addition, embodied in the principle of Occam's Razor, we often prefer stories that are the simplest; having the most support from our existing experiences.
This premise is explored in depth by \cite{storr2014unpersuadables}.

The use of narratives in science belies an understanding of natural phenomena through metaphor, the consequences of stories in science has been examined by \cite{mahoney2006tale,levy2008case,collier2011understanding,gelman2014stories}
Most recently this has been put forth to frame the decisions of economists in times of crisis by \cite{shriller2017narrative} and related to the political functions of democratic elections by

\todo{cite ``Multilingual Visual Sentiment Concept Clustering and Analysis'' paper}

%% Small piece: Can we determine if a character is bad?
%% Can we determine when a character dies?
%% Can we figure out the relationship between two characters?  Fighting?  A team? Romance?  And find when the relationship changes in nature?


Great defense on Friday! I was thinking about how else you could push on the shapes of story idea and am wondering whether you tried to predict the Library of Congress classification using the positivity time series. Maybe something like community detection after doing some kind of projection to deal with length differences? In my head, you would include both fiction and non-fiction and demonstrate that the six basic shapes uniquely identity fiction from most non-fiction.

Yeah I think we could feed the coefficient vector from the SVD for the first 10 or so modes (projection into that space) into a predictor and see how much predictive power we can get from each mode.
Clustering on those vectors would actually be straightfoward, do something like kMeans with k=2 and see if the two clusters are fiction/non-fiction?

There are many improvements possible for the visualizations hosted online at \href{http://teletherm.org/}{http://teletherm.org/}.
The teletherm animations can be improved through the use of the \verb|d3.timer| module for smoother animation.
Voronoi cells on the map are clipped at the boundary of the contiguous United States using a clipping mask that contains all 50 states as individual paths, and this does not work reliably in Google Chrome.
More issues for improvement are noted in the ``issues'' tab of the online source code repository at \href{https://github.com/andyreagan/teletherm.org}{https://github.com/andyreagan/teletherm.org}.
In addition, it will be possible to extend the teletherm project to incorporate temperature data from across the world.
